\section{The physical domain, the training targets and a fresh partition}\label{sec:domain}

The all-band scores of Table~\ref{tab:tails} mix three things: the error of the emulator, entries
where the retrieval is not defined, and entries where it is defined but badly conditioned. This
section separates them. It describes the physical domain of the retrieval and the entries of the
table that fall outside it, scores the main pipeline on that domain with its coverage and failures,
asks whether the out-of-domain entries matter when they are used as training targets, and repeats
the forward--inverse comparison on a fresh partition fixed in advance.

\subsection{The physical domain}
Write $t=Y_2+Y_3$, $s=Y_4$ and $q=1-\rho s$, with $\rho=0.7$. The retrieval analysis of
Section~\ref{sec:theory} holds on
\[
  \mathcal A=\{t>0,\ 0\le s<1\},
\]
which is a sufficient domain for the physical retrieval rather than the largest set on which the
rational formulas are defined: a negative albedo with $t>0$ gives $q>0$ and a well-defined inverse,
and $q<0$ violates the physical hypothesis without being a zero of a denominator. At $q=0$ the
forward formula is undefined, and at $t=0$, wherever it is defined, $L=Y_1$ does not depend on the
reflectance, so the band carries no information about it. A small positive transmission makes the
inverse ill-conditioned but not undefined.

\begin{table}[t]
\centering
\begin{tabular}{lrr}
\toprule
condition on the tabulated truth & entries & share \\
\midrule
transmitted flux $t=Y_2+Y_3 \le 0$ & 438 & $0.0066\%$ \\
$t < 10^{-12}$ & 32,130 & $0.4836\%$ \\
albedo $Y_4 < 0$ & 9,858 & $0.1484\%$ \\
albedo $Y_4 \ge 1$ & 454 & $0.0068\%$ \\
albedo exactly $2$ (fill marker) & 214 & $0.0032\%$ \\
\midrule
admissible $\mathcal{A}=\{t>0,\ 0\le Y_4<1\}$ & 6,633,849 & $99.8441\%$ \\
\bottomrule
\end{tabular}

\caption{The table against the physical domain, over all $23{,}313\times285=6{,}644{,}205$
state--band entries. The categories overlap. Small positive transmission is a conditioning
property, not a violation of $t>0$; an albedo outside $[0,1)$ violates the physical domain
without necessarily leaving the algebraic domain of the inverse.}
\label{tab:domain}
\end{table}

Table~\ref{tab:domain} counts the entries of the table outside $\mathcal A$. There are 438
nonpositive transmissions, the smallest $-4.5\times10^{-20}$, and $32{,}130$ transmissions below
$10^{-12}$; $9{,}858$ negative albedos, the smallest $-0.0724$; and 454 albedos at or above one,
214 of which equal $2$ exactly, which can only be a fill value. At 448 entries $q\le0$, the
smallest $-0.4$. In all, $10{,}356$ entries ($0.156\%$) lie outside $\mathcal A$. They occupy 227
bands but are concentrated near 1841--1908 and 2485--2493\,nm, and $18{,}859$ of the $23{,}313$
states are admissible at every band.

A stricter domain adds a conditioning floor,
\[
 \mathcal C=\mathcal A\cap\{t\ge\tau,\ q\ge\kappa\},
 \qquad \tau=10^{-12},\quad\kappa=0.3,
\]
with $\tau$ in the flux units of the table; at $\rho=0.7$ the albedo condition already gives
$q>0.3$. Inverting the exact tabulated components returns the reflectance to within
floating-point error, and that error is informative:

\begin{center}
\begin{tabular}{lrr}
\toprule
domain & round-trip RMSE (reflectance) & coverage \\
\midrule
all entries & $7.38\times10^{-3}$ & $100\%$ \\
$\mathcal A$ & $4.27\times10^{-3}$ & $99.844\%$ \\
$\mathcal C$ & $1.08\times10^{-9}$ & $99.516\%$ \\
\bottomrule
\end{tabular}
\end{center}

In exact arithmetic the round trip is exact wherever it is defined. In floating point, adding
$\rho t/q$ to $Y_1$ and subtracting $Y_1$ again loses the added term when it is small relative to
the spacing of numbers near $Y_1$, which is what happens at the smallest transmissions. Removing
$0.33\%$ of the entries takes the round-trip error from $4\times10^{-3}$ to $10^{-9}$; what remains
of the all-band reflectance tail on $\mathcal C$ belongs to the emulators.

\subsection{The main pipeline on the physical domain}\label{sec:tq-domain}
[[TQ-DOMAIN]]

\subsection{Training on physically admissible targets}\label{sec:tq-policy}
The training pipeline standardizes, compresses and fits the targets as they come, including the
$0.156\%$ of entries outside $\mathcal A$ and the fill values among them. Scoring on $\mathcal A$
cannot tell whether those entries shaped the fitted representations, the principal components or
the albedo plateau of Section~\ref{sec:scaling}; only refitting can. We therefore ran the pipeline
under three training policies: the original one (\emph{raw}); removal of every training state with
an entry outside $\mathcal A$ at any band (\emph{admissible}), with no clipping or imputation; and
an unfiltered random subset of the raw training rows of the same size as the admissible set
(\emph{matched}), which separates the effect of the filter from that of losing rows. Whole states
are removed so that every remaining spectrum is complete for the principal-component step; this
changes the training distribution, and the matched arm measures how much of any difference is
only the loss of $19\%$ of the rows. The validation rows, the validation loss and the test blocks
are those of the raw pipeline in all three arms, and every arm is scored on the same test entries
with a flux scale taken from the original training block.

The arms, the two network configurations of Table~\ref{tab:results} and Table~\ref{tab:wide}, the
seeds 101--110, the scoring and the summaries reported below were fixed in a written protocol
before the first run, and every run was completed and is reported. Each run records the digests of
its data and of its split indices, and saves the float64 predictions of every family.

[[TQ-POLICY]]

\subsection{Where the reflectance error sits}

\begin{figure}[p]
\centering
\includegraphics[height=0.80\textheight,width=\textwidth,keepaspectratio]{emit_quantiles_top_models.png}
\caption{Residual quantile bands across the test set of split 101, band by band, for six families of
the main pipeline. The envelopes are nested percentiles of the residual at each wavelength:
$25$--$75$ in red, $5$--$25$ and $75$--$95$ in blue, $2.5$--$5$ and $95$--$97.5$ in green,
$0.5$--$2.5$ and $97.5$--$99.5$ in grey, with the median as a thin black line. From the left: the
direct-flux residual $Y_2-\widehat{Y_2}$ in physical units; the radiance residual at $\rho=0.7$;
the reflectance residual over all bands; and the reflectance residual on the physical domain with
the floor $\tau_0=10^{-12}$. The last two columns share a vertical scale of $\pm2$ points.}
\label{fig:bands}
\end{figure}

Figure~\ref{fig:bands} shows where along the spectrum the errors of one split sit. The component
and radiance residuals are smooth in wavelength and widen where the transmitted flux is small, so
the families differ mainly in the amplitude of a shared shape. The reflectance residual does not
inherit that shape. Over all bands its outer envelopes open into vertical walls, bands where the
95th percentile of the absolute error exceeds two points, and the walls have two parts. A common
core of 54 bands, at 1126--1134, 1342--1469, 1796--1952 and 2412--2493\,nm, is a wall for every
family, the most accurate included: there the denominator, not the emulator, sets the width.
Outside the core each family adds bands in proportion to its inaccuracy: none for the kernel on
features, 3 for the stack, 12 for the input-scaled kernel, 22 for the corrected network, 29 for
the network and 79 for the cubic ridge, of a 133-band union. Restricting to the physical domain
changes none of these counts, and neither does the stricter floor $\tau_0=10^{-3}$, which removes
$6.3\%$ of the entries: the walls come from transmissions that are small but far above either
floor, where the retrieval divides a small component error by a small denominator. The feature
kernel's advantage in the tail is that it adds no wall of its own.

\subsection{A fresh partition, fixed in advance}\label{sec:confirm}
The comparisons above use partitions that were in view while the families were developed. To test
the reversal against that selection, a simplified pipeline was frozen together with one hypothesis
and evaluated once on a new partition. Its six families are the cubic ridge, a $3\times512$ network
(GELU activations, Adam) trained on the 285 standardized bands, the residual kernel of that
network, the input-scaled Mat\'ern kernel, the Mat\'ern kernel on the network's last hidden layer,
and a convex stack of the cubic, the corrected network, the input-scaled kernel and the feature
kernel. Every hyperparameter was copied from development on seed 101 and none was re-tuned. The
partition holds out 3{,}497 states, $15\%$ of the table, drawn with a seed used in no earlier run;
the remaining 19{,}816 states train the families, 1{,}982 of them serving only as the network's
early-stopping set. The held-out states belong to the same table as every other split, so this
evaluation guards against the choice of split and of hypothesis rather than testing
generalization to new states.

The hypothesis was that the family with the lowest radiance error does not have the lowest
all-band 95th-percentile reflectance error, and that its all-band 95th percentile is strictly larger
than that of the residual-corrected network.

\IfFileExists{table_confirmation.tex}{%
\begin{table}[t]
\centering
\resizebox{\textwidth}{!}{\begin{tabular}{lrrrrrr}
\toprule
family & radiance rel.\ $L_2$ [\%] & \multicolumn{3}{c}{all bands [pp]} & \multicolumn{2}{c}{admissible [pp]} \\
\cmidrule(lr){3-5}\cmidrule(lr){6-7}
 & & median & $p_{95}$ & $p_{99}$ & median & $p_{95}$ \\
\midrule
cubic ridge & 0.7544 & 0.475 & 70.72 & 129.1 & 0.468 & 70.67 \\
network $3\times512$ & 0.1719 & 0.106 & 41.83 & 84.1 & 0.105 & 28.07 \\
network + residual kernel & 0.0978 & 0.061 & 26.07 & 80.5 & 0.061 & 16.87 \\
input-scaled Mat\'ern kernel & 0.0900 & 0.028 & 26.89 & 189.4 & 0.027 & 19.38 \\
kernel on learned features & 0.0828 & 0.034 & 3.91 & 129.2 & 0.034 & 2.94 \\
convex stack & 0.0687 & 0.044 & 25.82 & 89.6 & 0.044 & 16.30 \\
\bottomrule
\end{tabular}
}
\caption{The fresh partition: 3{,}497 held-out states, evaluated once. The screened columns use
$t\ge10^{-12}$ and $q\ge0.3$, which covers $99.5327\%$ of the $996{,}645$ entries; no inversion is
non-finite in any family. The network here is the $3\times512$ GELU network of the frozen pipeline,
trained with Adam, not the network of Table~\ref{tab:results}.}
\label{tab:confirm}
\end{table}}{}

The reversal reproduces (Table~\ref{tab:confirm}). The convex stack has the lowest radiance error of
the six, $0.0687\%$, and the fourth-lowest reflectance tail; the kernel on learned features has a
larger radiance error, $0.0828\%$, and an all-band 95th percentile smaller by a factor of $6.60$, and
by $5.54$ on the screened entries. Both factors are close to the $6.13$ and $4.99$ of the same
pipeline on its development split, seed 101. The second clause of the hypothesis fails: the stack's
all-band 95th percentile is $25.82$ points and the corrected network's $26.07$, smaller by $0.25$
points where the hypothesis required it to be larger. The qualitative statement survives a single
read of a fresh partition, while the order among the four intermediate families does not separate at
this sample size.
